\section{Interpretation}
\label{sec:interpretation}

The comparisons identify three distinct sources of movement. Objective
weights determine which directions become active and at what parameter
scales. Ordered activation can make the central route longer than a
simultaneous motion toward an accurate endpoint. A feasible dual
completion can then require additional movement that is absent from the
primal projection. On the dyadic cube family these effects have separate,
matched asymptotic orders in both length and bounded local moves.

The exact target allocation is useful because a central endpoint need
not be the closest accurate point. It also shows why one aggregate rank
or determinant certificate may be insufficient: the entire objective
weight profile enters the optimum. Conversely, the sharp prefix factor
shows that ordered scalar activation imposes a strong restriction on
how inefficient a central route can be. For equal standard scales the
worst overhead is logarithmic in active rank under a square root, even
when the ambient matrix space is much larger.

Changing a barrier or a formulation requires a new metric comparison.
Optimal parameter alone does not eliminate the centrality cost, as the
canonical cube barriers and the explicit coupled radial family show.
A restriction of a cone barrier can also have much smaller primal
activity than the full homogeneous parameter suggests. The exposed-rank
and concrete-formulation results provide direct lower certificates when
auxiliary motion is permitted, rather than inferring primal cost from an
ambient parameter bound.

All these conclusions retain their stated contracts: fixed barrier
metric, specified starting point, actual terminal accuracy or feasible
duality gap, and bounded starting local norms. Within those contracts,
the explicit shortest routes and matching lower bounds make the cost of
centrality and the cost of dual completion separately measurable.
